
\documentclass{article}
\usepackage[utf8]{inputenc}
\usepackage{amsmath}
\usepackage{amsfonts}
\usepackage{amssymb}
\usepackage{amsthm}
\usepackage{boxedminipage}

\title{Master Manuscript V4: Structural Consistency and Residue Sanity}
\author{Thomas Hoffbauer}
\date{August 16, 2026}

\begin{document}

\maketitle

\section{The Residue Sanity Gate: Metric-Volumetric Invariance}

To ensure structural consistency in the EABC framework, we establish the Residue Sanity Gate. For any full-dimensional Euclidean lattice $\Lambda \subset \mathbb{R}^d$, the radial Epstein-Zeta function is defined as:
\[
Z_{\Lambda}^{\mathrm{len}}(s) = \sum_{0 \neq x \in \Lambda} \|x\|^{-s}.
\]
The fundamental geometric consistency condition at the pole $s=d$ is:
\begin{equation}
\boxed{\operatorname*{Res}_{s=d} Z_{\Lambda}^{\mathrm{len}}(s) = \frac{2\pi^{d/2}}{\Gamma(d/2)\operatorname{covol}(\Lambda)}}
\end{equation}

\noindent\textbf{Crucial Clarification on Invariance:} 
The residue is \emph{not} a topological invariant, but a \emph{metric-volumetric lattice invariant}. Under dilation $\Lambda \mapsto c\Lambda$, the zeta function transforms as $Z_{c\Lambda}^{\mathrm{len}}(s) = c^{-s}Z_\Lambda^{\mathrm{len}}(s)$, yielding:
\[
\operatorname*{Res}_{s=d} Z_{c\Lambda}^{\mathrm{len}} = c^{-d} \operatorname*{Res}_{s=d} Z_\Lambda^{\mathrm{len}}.
\]
More generally, for any $A \in GL(d, \mathbb{R})$, we have $\operatorname{covol}(A\Lambda) = |\det A|\operatorname{covol}(\Lambda)$, leading to:
\[
\boxed{
\operatorname*{Res}_{s=d} Z_{A\Lambda}^{\mathrm{len}} = |\det A|^{-1} \operatorname*{Res}_{s=d} Z_\Lambda^{\mathrm{len}}.
}
\]
Thus, the residue is invariant under all volume-preserving linear transformations ($|\det A| = 1$), including orientation-reversing ones.

\section{Divisor Information and Winding Audit (Argument Principle)}

For a meromorphic function, the winding number along a closed contour $\gamma$ in the resolvent set is given by the argument principle:
\[
\nu_\gamma = \frac{1}{2\pi i} \oint_\gamma \frac{Z'(s)}{Z(s)}\,ds = N_Z(\gamma) - N_P(\gamma) \in \mathbb{Z}.
\]
This size is invariant under homotopies of the contour within the complement of the divisor. This forms the clean dichotomy:
\[
\begin{array}{rcl}
\operatorname*{Res}_{s=d} Z_\Lambda &\longleftrightarrow& \text{metric/volumetric pole strength}, \\[2mm]
\displaystyle \frac{1}{2\pi i}\oint_\gamma d\log Z_\Lambda &\longleftrightarrow& \text{discrete divisor / winding information}.
\end{array}
\]

\section{Comparative Divisor and Winding Audit for $D_4$ and $E_8$}

Because $D_4$ lives in dimension 4 and $E_8$ in dimension 8, there is no canonical continuous cross-dimensional transition along which a single winding number could be transported directly from $D_4$ to $E_8$. Instead, we perform a rigorous comparative divisor audit by analyzing the divisors $\operatorname{div} Z_{D_4}^{\mathrm{len}}$ and $\operatorname{div} Z_{E_8}^{\mathrm{len}}$ using concentric contours centered at the respective dimensional poles:
\[
\gamma_{D_4}(R): |s-4|=R, \quad \gamma_{E_8}(R): |s-8|=R.
\]
The audit maps the discrete step functions:
\[
R \mapsto \bigl(\nu_{D_4}(R), \nu_{E_8}(R)\bigr),
\]
revealing the exact radial distances at which zeros and secondary poles appear and how internal factor singularities cancel out.

\section{Claim-Wall}

\begin{itemize}
    \item \textbf{Metric-Volumetric Residue Formula:} PROVEN
    \item \textbf{$E_8$- and $D_4$-Spezialisierungen:} PROVEN
    \item \textbf{Argument Principle / Winding Invariance:} PROVEN
    \item \textbf{Konkrete Produktdarstellung:} VALIDATED (only if pole structure matches the Gate)
    \item \textbf{Comparative Divisor Audit ($D_4$ vs $E_8$):} ESTABLISHED
    \item \textbf{Physikalische Interpretation:} HEURISTIC/NOT CLAIMED
\end{itemize}

end{document}
